% Options for packages loaded elsewhere
\PassOptionsToPackage{unicode}{hyperref}
\PassOptionsToPackage{hyphens}{url}
\PassOptionsToPackage{dvipsnames,svgnames,x11names}{xcolor}
%
\documentclass[
  11pt,
  a4paper,
]{article}
\usepackage{amsmath,amssymb}
\usepackage{iftex}
\ifPDFTeX
  \usepackage[T1]{fontenc}
  \usepackage[utf8]{inputenc}
  \usepackage{textcomp} % provide euro and other symbols
\else % if luatex or xetex
  \usepackage{unicode-math} % this also loads fontspec
  \defaultfontfeatures{Scale=MatchLowercase}
  \defaultfontfeatures[\rmfamily]{Ligatures=TeX,Scale=1}
\fi
\usepackage{lmodern}
\ifPDFTeX\else
  % xetex/luatex font selection
\fi
% Use upquote if available, for straight quotes in verbatim environments
\IfFileExists{upquote.sty}{\usepackage{upquote}}{}
\IfFileExists{microtype.sty}{% use microtype if available
  \usepackage[]{microtype}
  \UseMicrotypeSet[protrusion]{basicmath} % disable protrusion for tt fonts
}{}
\makeatletter
\@ifundefined{KOMAClassName}{% if non-KOMA class
  \IfFileExists{parskip.sty}{%
    \usepackage{parskip}
  }{% else
    \setlength{\parindent}{0pt}
    \setlength{\parskip}{6pt plus 2pt minus 1pt}}
}{% if KOMA class
  \KOMAoptions{parskip=half}}
\makeatother
\usepackage{xcolor}
\usepackage{longtable,booktabs,array}
\usepackage{calc} % for calculating minipage widths
% Correct order of tables after \paragraph or \subparagraph
\usepackage{etoolbox}
\makeatletter
\patchcmd\longtable{\par}{\if@noskipsec\mbox{}\fi\par}{}{}
\makeatother
% Allow footnotes in longtable head/foot
\IfFileExists{footnotehyper.sty}{\usepackage{footnotehyper}}{\usepackage{footnote}}
\makesavenoteenv{longtable}
\setlength{\emergencystretch}{3em} % prevent overfull lines
\providecommand{\tightlist}{%
  \setlength{\itemsep}{0pt}\setlength{\parskip}{0pt}}
\setcounter{secnumdepth}{-\maxdimen} % remove section numbering
% Layout of the GZero documents (used by docs/build_manual.py through pandoc).
\usepackage[a4paper,top=25mm,bottom=25mm,left=25mm,right=25mm,headheight=14pt]{geometry}
\usepackage{lmodern}
\renewcommand{\familydefault}{\sfdefault}
\usepackage[T1]{fontenc}
\usepackage{microtype}
\usepackage{amsmath,amssymb}
\usepackage{xcolor}
\definecolor{gzblue}{HTML}{0072B2}
\definecolor{gzgrey}{HTML}{F3F4F6}
\usepackage{titlesec}
\titleformat{\section}{\Large\bfseries\color{gzblue}}{}{0pt}{}[\vspace{2pt}\titlerule]
\titleformat{\subsection}{\large\bfseries}{}{0pt}{}
\titleformat{\subsubsection}{\normalsize\bfseries}{}{0pt}{}
\titlespacing*{\section}{0pt}{18pt plus 4pt minus 2pt}{8pt}
\titlespacing*{\subsection}{0pt}{12pt plus 3pt minus 2pt}{4pt}
\usepackage{fancyhdr}
\pagestyle{fancy}
\fancyhf{}
\makeatletter\AtBeginDocument{\fancyhead[L]{\small\@title}}\makeatother
\fancyhead[R]{\small\nouppercase{\rightmark}}
\fancyfoot[C]{\small\thepage}
\renewcommand{\headrulewidth}{0.4pt}
\renewcommand{\sectionmark}[1]{\markboth{#1}{#1}}
\renewcommand{\subsectionmark}[1]{}
\fancypagestyle{plain}{\fancyhf{}\fancyfoot[C]{\small\thepage}\renewcommand{\headrulewidth}{0pt}}
\setlength{\parindent}{0pt}
\setlength{\parskip}{5pt plus 1pt}
\usepackage{etoolbox}
% tables: smaller type, a little more row height, ragged cells
\AtBeginEnvironment{longtable}{\small\renewcommand{\arraystretch}{1.25}}
\usepackage{ragged2e}
\setlength{\LTpre}{6pt}
\setlength{\LTpost}{6pt}
\usepackage{enumitem}
\setlist{itemsep=2pt,topsep=4pt}
% code blocks on a light background
\usepackage{fvextra}
\DefineVerbatimEnvironment{Highlighting}{Verbatim}{breaklines,commandchars=\\\{\}}
\BeforeBeginEnvironment{verbatim}{\begingroup\small}
\AfterEndEnvironment{verbatim}{\endgroup}
\usepackage[most]{tcolorbox}
\usepackage{tocloft}
\renewcommand{\cftsecfont}{\bfseries}
\renewcommand{\cftsecpagefont}{\bfseries}
\setlength{\cftbeforesecskip}{3pt}
\setcounter{tocdepth}{2}
\usepackage{hyperref}
\hypersetup{colorlinks=true,linkcolor=gzblue,urlcolor=gzblue,citecolor=gzblue,pdfauthor={arman-rashid}}
\urlstyle{same}
\sloppy
\ifLuaTeX
  \usepackage{selnolig}  % disable illegal ligatures
\fi
\IfFileExists{bookmark.sty}{\usepackage{bookmark}}{\usepackage{hyperref}}
\IfFileExists{xurl.sty}{\usepackage{xurl}}{} % add URL line breaks if available
\urlstyle{same}
\hypersetup{
  pdftitle={GZero 1.3 and 1.4: new analysis methods},
  colorlinks=true,
  linkcolor={Maroon},
  filecolor={Maroon},
  citecolor={Blue},
  urlcolor={Blue},
  pdfcreator={LaTeX via pandoc}}

\title{GZero 1.3 and 1.4: new analysis methods}
\usepackage{etoolbox}
\makeatletter
\providecommand{\subtitle}[1]{% add subtitle to \maketitle
  \apptocmd{\@title}{\par {\large #1 \par}}{}{}
}
\makeatother
\subtitle{October 2026, arman-rashid}
\author{}
\date{}

\begin{document}
\maketitle

GZero 1.3 adds three methods that are standard in break-junction papers:
the junction yield with per-trace plateau statistics, single-level model
fits of I-V curves, and transition voltage spectroscopy. GZero 1.4 adds
the analysis methods of Rashid et al., JACS 2025 (finite-temperature
energy-grid fit, noise protocol, event detection, dwell time) and an
EGaIn ensemble tab built after the Hong (XMe) and Chiechi (GaussFit)
codes.

\hypertarget{whats-new}{%
\section{What's new}\label{whats-new}}

Each new method answers a question the existing tabs couldn't. All of
them run on data GZero already loads, so no new file formats are needed.

\begin{longtable}[]{@{}
  >{\raggedright\arraybackslash}p{(\columnwidth - 6\tabcolsep) * \real{0.2124}}
  >{\raggedright\arraybackslash}p{(\columnwidth - 6\tabcolsep) * \real{0.2019}}
  >{\raggedright\arraybackslash}p{(\columnwidth - 6\tabcolsep) * \real{0.3738}}
  >{\raggedright\arraybackslash}p{(\columnwidth - 6\tabcolsep) * \real{0.1820}}@{}}
\toprule\noalign{}
\begin{minipage}[b]{\linewidth}\raggedright
Addition
\end{minipage} & \begin{minipage}[b]{\linewidth}\raggedright
Where
\end{minipage} & \begin{minipage}[b]{\linewidth}\raggedright
Question it answers
\end{minipage} & \begin{minipage}[b]{\linewidth}\raggedright
Module
\end{minipage} \\
\midrule\noalign{}
\endhead
\bottomrule\noalign{}
\endlastfoot
Junction statistics & New tab after Plateau length & How many traces
actually had a molecule, at what conductance, and did that drift during
the run? & \texttt{gzero/\allowbreak{}junctions.py} \\
Single-level model fit & I-V tab & Where is the molecular level
($\varepsilon_0$) and how strongly is it coupled to the electrodes
($\Gamma$)? & \texttt{gzero/\allowbreak{}ivmodels.py} \\
Transition voltage spectroscopy & I-V tab & Where does the
Fowler-Nordheim plot turn over, and what level position does that imply?
& \texttt{gzero/\allowbreak{}ivmodels.py} \\
Events & New tab & Which traces flicker or switch mechanically? &
\texttt{gzero/\allowbreak{}events.py} \\
EGaIn (ensemble) & New tab & Statistics of large-area junctions, yield,
rectification and $\beta$ & \texttt{gzero/\allowbreak{}egain.py} \\
\end{longtable}

One existing function got a fix along the way. The averaged I-V curve
used to take the overlap of all sweeps, so the half sweeps at the start
and end of a recording could shrink it to nothing. It now averages only
sweeps that cover the typical bias range.

\hypertarget{papers-reviewed}{%
\section{Papers reviewed}\label{papers-reviewed}}

I went through the methods sections of commonly cited break-junction
papers and looked for analyses people report routinely that GZero didn't
have yet. Four papers define what was added in 1.3:

\begin{longtable}[]{@{}
  >{\raggedright\arraybackslash}p{(\columnwidth - 4\tabcolsep) * \real{0.2116}}
  >{\raggedright\arraybackslash}p{(\columnwidth - 4\tabcolsep) * \real{0.3979}}
  >{\raggedright\arraybackslash}p{(\columnwidth - 4\tabcolsep) * \real{0.3605}}@{}}
\toprule\noalign{}
\begin{minipage}[b]{\linewidth}\raggedright
Paper
\end{minipage} & \begin{minipage}[b]{\linewidth}\raggedright
What they do
\end{minipage} & \begin{minipage}[b]{\linewidth}\raggedright
What GZero does with it
\end{minipage} \\
\midrule\noalign{}
\endhead
\bottomrule\noalign{}
\endlastfoot
\href{https://doi.org/10.1103/PhysRevLett.102.126803}{Kamenetska et al.,
PRL 102, 126803 (2009)} & Count how many breaking traces show a
molecular plateau (junction formation probability) and how long the
plateaus are, for a series of molecules & Junction yield per trace with
a binomial error, plus yield and conductance in blocks of traces over
the run \\
\href{https://doi.org/10.1002/smll.200902227}{Zotti et al., Small 6,
1529 (2010)} & Fit MCBJ I-V curves of tolanes with thiol, nitro and
cyano anchors with a single-level model to get $\varepsilon_0$ and
$\Gamma$ & Single-level fit of every curve and of the mean forward and
backward curves, with errors and an optional bias asymmetry \\
\href{https://doi.org/10.1103/PhysRevLett.97.026801}{Beebe et al., PRL
97, 026801 (2006)} & Plot ln(I/$V^{2}$) against 1/V and read off the
minimum, the transition voltage $V_t$; it scales with the HOMO offset
from photoemission & $V_t$ for both polarities of every curve, with a
Fowler-Nordheim plot of the mean curves \\
\href{https://doi.org/10.1103/PhysRevB.85.035442}{Baldea, PRB 85, 035442
(2012)} & Shows that both polarities of $V_t$ together give the level
position $\varepsilon_0$ and the bias asymmetry in closed form &
$\varepsilon_0$ and the asymmetry a from $V_t^{+}$ and $V_t^{-}$, so TVS
and the fit can be compared curve by curve \\
\end{longtable}

The 2D correlation, clustering and flicker-noise methods from earlier
papers were already in GZero. I left out methods that need hardware most
setups don't have, like thermopower and shot noise.

\hypertarget{junction-statistics}{%
\section{Junction statistics}\label{junction-statistics}}

A trace counts as a molecular junction if it stays flat inside the
molecular window for at least 0.1 nm. The yield is the fraction of
traces that do.

The tricky part is ``flat''. A trace that only tunnels still passes
through a two-decade window in about 0.1-0.2 nm, so a plain
length-in-window test counts almost every trace. GZero therefore checks
the local slope at each point:

\[
s(z) = \frac{|\log G(z + w/2) - \log G(z - w/2)|}{w}
\]

w is 0.05 nm by default. Tunnelling falls at $\beta$ / ln 10 decades per
nm: about 9 for gold in vacuum and 4-5 in solvents. Molecular plateaus
usually fall by well under 1 decade per nm. A point counts as flat below
2 decades/nm, which sits between the two.

Per trace you get three numbers:

\begin{itemize}
\tightlist
\item
  \textbf{flat length}: summed distance of the flat points inside the
  window
\item
  \textbf{plateau conductance}: median log G of those points
\item
  \textbf{junction yes/no}: flat length of at least \emph{min plateau}
\end{itemize}

From these the tab gives the yield p with its binomial error
sqrt(p(1-p)/N), the median plateau conductance and length, and a drift
in decades per 1000 traces. The drift is a straight line fitted to the
block medians, with 100 traces per block by default. The five plots are
the plateau conductance histogram, conductance against trace number,
yield against trace number, the flat-length histogram with the threshold
marked, and length against conductance.

Use peak from histogram fit sets the window to the fitted peak +/- 3
$\sigma$. Then look at the flat-length histogram. Tunnelling traces pile
up near zero and molecular ones form a second hump, and the threshold
belongs in the dip between them.

\hypertarget{single-level-model}{%
\section{Single-level model}\label{single-level-model}}

The fit describes each I-V curve with one molecular level at energy
$\varepsilon_0$ from the Fermi level, broadened by the coupling $\Gamma$
to each electrode. At zero temperature the Landauer current through a
Lorentzian transmission has a closed form:

\[
I(V) = N G_0 \Gamma \left[\arctan\frac{V/2 - \varepsilon(V)}{\Gamma} + \arctan\frac{V/2 + \varepsilon(V)}{\Gamma}\right], \quad \varepsilon(V) = \varepsilon_0 + aV
\]

Energies are in eV and V in volts, so $G_0$ times an energy in eV is
directly a current in amperes. Both electrodes are coupled equally. N is
the number of molecules in parallel and stays at 1. a lets the level
move with the bias in an asymmetric junction; it is 0 unless you tick
\emph{fit asymmetry a}. At low bias the model gives G/$G_0$ = $\Gamma^2$
/ ($\varepsilon_0^2$ + $\Gamma^2$).

How the fit works:

\begin{enumerate}
\def\labelenumi{\arabic{enumi}.}
\tightlist
\item
  The current is scaled by its largest value so all curves weigh the
  same.
\item
  A least-squares fit starts from four guesses of $\varepsilon_0$ (0.4,
  0.8, 1.5 and 3 times the largest bias). Each guess gets a $\Gamma$
  that matches the measured low-bias conductance.
\item
  The fit with the smallest residual is kept. Errors come from the fit
  covariance, and $R^{2}$ is reported.
\item
  This is done for every curve (up to 500), for the mean forward and
  backward curves, and for the most probable curve.
\end{enumerate}

The model can't tell HOMO from LUMO transport, because both give the
same I-V. $\varepsilon_0$ is therefore reported as a positive number.

\hypertarget{transition-voltage-spectroscopy}{%
\section{Transition voltage
spectroscopy}\label{transition-voltage-spectroscopy}}

The transition voltage $V_t$ is where $V^{2}$/|I| peaks, which is the
same as the minimum of the Fowler-Nordheim plot, ln(|I|/$V^{2}$) against
1/V. It needs no fit, which is why so many papers report it.

For each curve and polarity GZero averages the current in 60 bias bins
above 0.05 V and finds the largest $V^{2}$/|I|. It then refines that
point with a parabola through it and its two neighbours. If the maximum
sits at the end of the sweep, the bias never reached $V_t$, and the
value is left empty rather than guessed.

In the single-level model with $\Gamma$ much smaller than
$\varepsilon_0$, the two polarities give the level position and the
asymmetry exactly (Baldea 2012):

\[
\varepsilon_0 = \frac{2|V_t^+ V_t^-|}{\sqrt{(V_t^+)^2 + (V_t^-)^2 + \frac{10}{3}|V_t^+ V_t^-|}}, \quad a = \frac{(V_t^+ + V_t^-)\,\varepsilon_0}{4|V_t^+ V_t^-|}
\]

For a symmetric junction this reduces to $\varepsilon_0$ = (sqrt(3)/2)
$V_t$, about 0.87 $V_t$. The a formula follows from the same
stationarity condition and was checked numerically against the exact
model (next section). Comparing the $\varepsilon_0$ from TVS with the
$\varepsilon_0$ from the fit is a quick test of whether a single level
describes the curves at all.

\hypertarget{validation-of-the-1.3-methods}{%
\section{Validation of the 1.3
methods}\label{validation-of-the-1.3-methods}}

On data with known answers, every new method recovers the true value
within its error. The SLM and TVS rows are model curves with 2 \% noise.
The yield rows are GZero's own simulator going through the normal trace
detection.

\begin{longtable}[]{@{}
  >{\raggedright\arraybackslash}p{(\columnwidth - 4\tabcolsep) * \real{0.2837}}
  >{\raggedright\arraybackslash}p{(\columnwidth - 4\tabcolsep) * \real{0.2891}}
  >{\raggedright\arraybackslash}p{(\columnwidth - 4\tabcolsep) * \real{0.3973}}@{}}
\toprule\noalign{}
\begin{minipage}[b]{\linewidth}\raggedright
Test
\end{minipage} & \begin{minipage}[b]{\linewidth}\raggedright
True value
\end{minipage} & \begin{minipage}[b]{\linewidth}\raggedright
GZero
\end{minipage} \\
\midrule\noalign{}
\endhead
\bottomrule\noalign{}
\endlastfoot
SLM fit, symmetric & $\varepsilon_0$ 0.800 eV, $\Gamma$ 10 meV & 0.801
eV, 10.0 meV ($R^{2}$ 0.9997) \\
SLM fit, asymmetric & $\varepsilon_0$ 0.800 eV, a 0.100 & 0.795 eV, a
0.096 \\
TVS, symmetric & $\varepsilon_0$ 0.800 eV & $V_t$ = +/-0.924 V,
$\varepsilon_0$ 0.800 eV \\
TVS, asymmetric & $\varepsilon_0$ 0.800 eV, a 0.100 & $V_t^{+}$ 1.191 V,
$V_t^{-}$ -0.746 V, $\varepsilon_0$ 0.800 eV, a 0.100 \\
TVS, asymmetric & $\varepsilon_0$ 1.000 eV, a -0.080 & $\varepsilon_0$
1.002 eV, a -0.080 \\
Self-test I-V & $\varepsilon_0$ 0.800 eV, $\Gamma$ 20 meV & fit 0.800
eV, 20.0 meV; TVS 0.801 eV \\
Junction yield, 200 traces & 70 \% & 68.5 +/- 3.3 \% \\
Junction yield, 200 traces & 30 \% & 29.0 +/- 3.2 \% \\
Plateau conductance & log G -3.5 & median -3.55 \\
\end{longtable}

The slight shift of the plateau conductance comes from the simulated
plateau slope (0.3 decades/nm), which pulls the median a little below
the plateau's start.

\hypertarget{version-1.4-methods-from-rashid-et-al.-2025}{%
\section{Version 1.4: methods from Rashid et
al.~2025}\label{version-1.4-methods-from-rashid-et-al.-2025}}

Every analysis section of the SI of
\href{https://doi.org/10.1021/jacs.4c13531}{Rashid et al., JACS 147, 830
(2025)} is now in GZero, except the DFT and transport calculations.

\begin{longtable}[]{@{}
  >{\raggedright\arraybackslash}p{(\columnwidth - 4\tabcolsep) * \real{0.1614}}
  >{\raggedright\arraybackslash}p{(\columnwidth - 4\tabcolsep) * \real{0.4462}}
  >{\raggedright\arraybackslash}p{(\columnwidth - 4\tabcolsep) * \real{0.3624}}@{}}
\toprule\noalign{}
\begin{minipage}[b]{\linewidth}\raggedright
SI section
\end{minipage} & \begin{minipage}[b]{\linewidth}\raggedright
What the paper does
\end{minipage} & \begin{minipage}[b]{\linewidth}\raggedright
In GZero
\end{minipage} \\
\midrule\noalign{}
\endhead
\bottomrule\noalign{}
\endlastfoot
2.3 Flicker noise & Hold 160 ms at 100 mV, drop 10 ms, keep holds whose
first and last 100 points sit within +/- 1 $\sigma$ of the peak, DFT
squared, 100 Hz-1 kHz, n where Pearson r is smallest (0.3-2.3, step
0.01) & Flicker noise tab, preset ``Rashid et al.~2025''; new \emph{end
points} setting; minimum-\textbar r\textbar{} grid value reported \\
2.4 I-V fit & Landauer integral with Fermi functions on an energy grid
(-10 to 10 eV, 0.02 meV), $\Gamma_L$ = $\Gamma_R$, Levenberg-Marquardt
fit of the most probable I-V & I-V tab, preset ``Rashid et al.~2025'';
energy grid or exact digamma form; most probable curve from a Gaussian
per bias bin \\
2.5 Mechanical modulation & 2D conductance-time histograms of holds with
the most probable G per time & Events tab with segments = holds \\
2.6 Event detection & Butterworth 200 Hz, Savitzky-Golay (order 20, 20
side points), derivative, peak thresholds 0.0007 / 0.002, width 50 &
Events tab \\
2.7 Dwell time & $\tau$ = A/w x 1000/f\_s ms from the peak area &
Histograms tab, with every peak fit \\
\end{longtable}

At finite temperature the Landauer integral is evaluated either on the
energy grid of the paper or with the exact closed form, which uses the
digamma function psi:

\[
\int T(E)\, f(E - \mu)\, dE = \pi\Gamma \left[\frac{1}{2} - \frac{1}{\pi}\, \mathrm{Im}\, \psi\!\left(\frac{1}{2} + \frac{\Gamma - i(\mu - \varepsilon)}{2\pi k_B T}\right)\right]
\]

The two agree to $10^{-9}$. Fits therefore start from the closed form
and the grid only polishes, which brings a grid fit down from 50 s to
under 1 s.

\begin{longtable}[]{@{}
  >{\raggedright\arraybackslash}p{(\columnwidth - 4\tabcolsep) * \real{0.4309}}
  >{\raggedright\arraybackslash}p{(\columnwidth - 4\tabcolsep) * \real{0.2745}}
  >{\raggedright\arraybackslash}p{(\columnwidth - 4\tabcolsep) * \real{0.2645}}@{}}
\toprule\noalign{}
\begin{minipage}[b]{\linewidth}\raggedright
Test
\end{minipage} & \begin{minipage}[b]{\linewidth}\raggedright
True
\end{minipage} & \begin{minipage}[b]{\linewidth}\raggedright
GZero 1.4
\end{minipage} \\
\midrule\noalign{}
\endhead
\bottomrule\noalign{}
\endlastfoot
Preset fit, model curve with the G1 values of Table S1, 300 K &
$\varepsilon_0$ 1.10 eV, $\Gamma$ 11.3 meV & 1.100 eV, 11.30 meV \\
Same, 2 \% noise, through the GUI (mean, most probable, single curves) &
1.10 eV, 11.3 meV & 1.098-1.100 eV, 11.3 meV \\
Noise preset, 5 simulated sets, holds on the plateau & n = 1.5 & 1.53
(Morris preset: 1.59) \\
Dwell time, 70 \% molecular, 0.4 nm at 2 nm/s & 140 ms & 135-146 ms \\
\end{longtable}

The SI gives no temperature, so the preset assumes 300 K. Change it if
your measurements ran at another temperature.

\hypertarget{events-tab}{%
\section{Events tab}\label{events-tab}}

The tab finds flickering and mechanical switching in G(t) with the four
steps of SI 2.6, and classifies each trace as mechanical, flickering or
quiet.

\begin{enumerate}
\def\labelenumi{\arabic{enumi}.}
\tightlist
\item
  Butterworth low-pass, 200 Hz, order 2, applied forwards and backwards
  so events aren't shifted in time.
\item
  Savitzky-Golay smoothing, order 20 over 41 points. scipy's own
  coefficients collapse to about $10^{-17}$ at this order, so GZero
  builds the filter from a Legendre-basis QR. It matches scipy to
  $10^{-14}$ at low order.
\item
  Second-order central derivative dG/dt with G in $G_0$.
\item
  Peaks above 0.0007 are flickers and above 0.002 mechanical events.
  Events outside the molecular log G range are dropped as false
  positives.
\end{enumerate}

The plots are an example trace with the events marked, its derivative
with both thresholds, the share of each trace class, flickers per trace,
and a before/after transition map. With holds you also get the
conductance-time histogram with the most probable G.

\textbf{Three choices to confirm.} The SI doesn't pin these down, so
GZero copies LabVIEW's defaults:

\begin{itemize}
\tightlist
\item[$\square$]
  \textbf{dt = 20 ms is LabVIEW's dt input}, so dG/dt = (G{[}i+1{]} -
  G{[}i-1{]}) / (2 x 0.02 s). The thresholds then correspond to steps of
  about 3 x $10^{-4}$ $G_0$ (flicker) and 9 x $10^{-4}$ $G_0$
  (mechanical), since the dG/dt peak is about 2.2 times the step.
\item[$\square$]
  \textbf{``Width 50'' is the Peak Detector's quadratic-fit window}, not
  a minimum peak width. Read as a minimum width at half height, it
  rejects almost every step, because 200 Hz-filtered steps give
  derivative peaks only 20-30 samples wide at 10 kHz. Both readings are
  selectable.
\item[$\square$]
  \textbf{Only rises are counted} (Peak Detector default), so a flicker
  up and back counts once. \emph{Count = both} counts every step.
\end{itemize}

On synthetic steps a 4 x $10^{-4}$ $G_0$ flicker at $10^{-3}$ $G_0$ is
found as a flicker, a jump from $10^{-6}$ to 1.6 x $10^{-3}$ $G_0$ as
mechanical, and a flat trace with 3 \% noise gives nothing.

\hypertarget{egain-ensemble-tab}{%
\section{EGaIn ensemble tab}\label{egain-ensemble-tab}}

Large-area junctions are analysed per junction, not per sweep. 20
junctions with 20 sweeps each are 20 independent measurements, not 400.
The tab follows \href{https://doi.org/10.1021/jp210445y}{Reus et
al.~2012}, as XMe and GaussFit also do.

\begin{enumerate}
\def\labelenumi{\arabic{enumi}.}
\tightlist
\item
  Each file is one junction. J = I / A, with A from the contact diameter
  (XMe's approach) or typed in.
\item
  A junction whose \textbar J\textbar{} reaches the short threshold (100
  A/cm2) is a short. The yield is working junctions over all junctions.
\item
  Sweeps go onto a common bias grid. At each bias,
  log10\textbar J\textbar{} is histogrammed and fitted with a Gaussian
  to give the centre and $\sigma_\mathrm{log}$.
\item
  The confidence interval uses the number of junctions $n_j$ as degrees
  of freedom (GaussFit's convention); formula below.
\item
  Also computed: log R = log\textbar J(+V)\textbar{} -
  log\textbar J(-V)\textbar{} per cycle, $V_\mathrm{trans}$ per cycle,
  and a single-level fit of the mean J(V) with free N.
\item
  Fit $\beta$ across data sets with typed-in lengths d (formula below).
  Each point is weighted by the standard error of its mean
  (\href{https://doi.org/10.1021/ja408652h}{Simeone et al.~2013} report
  0.92 per carbon).
\end{enumerate}

\[
\mathrm{CI} = t_{1-\alpha/2,\; n_j - 1}\, \frac{\sigma_\mathrm{log}}{\sqrt{n_j - 1}}, \qquad \log_{10} |J| = \log_{10} J_0 - \frac{\beta\, d}{\ln 10}
\]

\begin{longtable}[]{@{}
  >{\raggedright\arraybackslash}p{(\columnwidth - 4\tabcolsep) * \real{0.3812}}
  >{\raggedright\arraybackslash}p{(\columnwidth - 4\tabcolsep) * \real{0.2984}}
  >{\raggedright\arraybackslash}p{(\columnwidth - 4\tabcolsep) * \real{0.2904}}@{}}
\toprule\noalign{}
\begin{minipage}[b]{\linewidth}\raggedright
Check (simulated junctions)
\end{minipage} & \begin{minipage}[b]{\linewidth}\raggedright
True
\end{minipage} & \begin{minipage}[b]{\linewidth}\raggedright
Found
\end{minipage} \\
\midrule\noalign{}
\endhead
\bottomrule\noalign{}
\endlastfoot
$\beta$, 5 sets of 30 junctions & 0.9 & 0.904 +/- 0.016 \\
Rectification log R at 0.5 V & 0.30 & 0.30 +/- 0.08 \\
Mean log\textbar J\textbar{} per bias & as simulated & inside its 95 \%
CI \\
$\beta$ error coverage, 20 repeats & 68 \% within 1 $\sigma$ & 65 \% \\
\end{longtable}

The first version weighted $\beta$ by $\sigma_\mathrm{log}$ and claimed
errors 3 times too small. Weighting by the standard error fixed that,
which is what the coverage row shows.

\hypertarget{open-source-code-reviewed}{%
\section{Open-source code reviewed}\label{open-source-code-reviewed}}

Two public codes cover the same ground and were read in full for their
EGaIn and I-V methods.

\begin{longtable}[]{@{}
  >{\raggedright\arraybackslash}p{(\columnwidth - 8\tabcolsep) * \real{0.2360}}
  >{\raggedright\arraybackslash}p{(\columnwidth - 8\tabcolsep) * \real{0.1550}}
  >{\raggedright\arraybackslash}p{(\columnwidth - 8\tabcolsep) * \real{0.1145}}
  >{\raggedright\arraybackslash}p{(\columnwidth - 8\tabcolsep) * \real{0.2223}}
  >{\raggedright\arraybackslash}p{(\columnwidth - 8\tabcolsep) * \real{0.2422}}@{}}
\toprule\noalign{}
\begin{minipage}[b]{\linewidth}\raggedright
Code
\end{minipage} & \begin{minipage}[b]{\linewidth}\raggedright
Group
\end{minipage} & \begin{minipage}[b]{\linewidth}\raggedright
License
\end{minipage} & \begin{minipage}[b]{\linewidth}\raggedright
Modules
\end{minipage} & \begin{minipage}[b]{\linewidth}\raggedright
How GZero used it
\end{minipage} \\
\midrule\noalign{}
\endhead
\bottomrule\noalign{}
\endlastfoot
\href{https://github.com/Pilab-XMU/XMe_DataAnalysis}{XMe\_DataAnalysis}
& Wenjing Hong, Xiamen University & Apache 2.0 & Basic histograms,
correlation, spectral clustering, PSD noise, I-V, single trace,
thermopower, electrochemistry, EGaIn & EGaIn: J from contact diameter,
log\textbar J\textbar{} per bias with Gaussian fits. The methods were
re-implemented; XMe asks to be cited when compared with \\
\href{https://github.com/rchiechi/gaussfit}{GaussFit} & Ryan Chiechi &
GPL-3.0 & EGaIn / CP-AFM parsing, Gaussian log\textbar J\textbar,
rectification, $V_\mathrm{trans}$, dJ/dV, lag filter, SLM &
Junction-based CI, rectification and $V_\mathrm{trans}$ per cycle.
Methods only: GPL code can't go into Apache-licensed GZero \\
\end{longtable}

GZero already covers the other Hong-group analyses in XMe: correlation,
spectral clustering and flicker noise. Thermopower and electrochemical
gating need measurement types GZero doesn't read yet.

\hypertarget{sources}{%
\section{Sources}\label{sources}}

Code: \href{https://github.com/arman-rashid/GZero}{arman-rashid/GZero},
release v1.4. The full method description is in the user manual
(\texttt{docs/\allowbreak{}User\_\allowbreak{}Manual.pdf}), sections
5.2, 6.4, 9.5, 9.8 and 10.2-10.5, references 26-35.

\begin{itemize}
\tightlist
\item
  M. Kamenetska et al., ``Formation and evolution of single-molecule
  junctions'', Phys. Rev.~Lett. 102, 126803 (2009).
  \href{https://doi.org/10.1103/PhysRevLett.102.126803}{doi:10.1103/PhysRevLett.102.126803}
\item
  L. A. Zotti et al., ``Revealing the role of anchoring groups in the
  electrical conduction through single-molecule junctions'', Small 6,
  1529-1535 (2010).
  \href{https://doi.org/10.1002/smll.200902227}{doi:10.1002/smll.200902227}
\item
  J. M. Beebe et al., ``Transition from direct tunneling to field
  emission in metal-molecule-metal junctions'', Phys. Rev.~Lett. 97,
  026801 (2006).
  \href{https://doi.org/10.1103/PhysRevLett.97.026801}{doi:10.1103/PhysRevLett.97.026801}
\item
  I. Baldea, ``Ambipolar transition voltage spectroscopy: analytical
  results and experimental agreement'', Phys. Rev.~B 85, 035442 (2012).
  \href{https://doi.org/10.1103/PhysRevB.85.035442}{doi:10.1103/PhysRevB.85.035442}
\item
  U. Rashid et al., ``Mapping the extended ground state reactivity
  landscape of a photoswitchable molecule at a single molecular level'',
  J. Am. Chem. Soc. 147, 830-840 (2025).
  \href{https://doi.org/10.1021/jacs.4c13531}{doi:10.1021/jacs.4c13531}
  (SI sections 2.3-2.7)
\item
  L. J. Reus et al., ``Statistical tools for analyzing measurements of
  charge transport'', J. Phys. Chem. C 116, 6714 (2012).
  \href{https://doi.org/10.1021/jp210445y}{doi:10.1021/jp210445y}
\item
  F. C. Simeone et al., ``Defining the value of injection current and
  effective electrical contact area for EGaIn-based molecular tunneling
  junctions'', J. Am. Chem. Soc. 135, 18131 (2013).
  \href{https://doi.org/10.1021/ja408652h}{doi:10.1021/ja408652h}
\item
  E. M. Opodi et al., ``A single level tunneling model for molecular
  junctions: evaluating the simulation methods'', Phys. Chem. Chem.
  Phys. 24, 11958 (2022).
  \href{https://doi.org/10.1039/D1CP05807J}{doi:10.1039/D1CP05807J}
\item
  \href{https://github.com/Pilab-XMU/XMe_DataAnalysis}{XMe\_DataAnalysis}
  and \href{https://github.com/rchiechi/gaussfit}{GaussFit} on GitHub
\end{itemize}

\end{document}
